\section{Words and effective language algebra}
\label{sec:words}

\subsection{A numerical free monoid}

Fix $q\ge1$ and write $\Sigma_q=\{0,\ldots,q-1\}$. Put $b=q+1$.
Encode a word by
\[
 \operatorname{code}(\eps)=0,\qquad
 \operatorname{code}(wa)=b\operatorname{code}(w)+(a+1).
\]
Thus a valid nonzero code has only digits $1,\ldots,q$ in base $b$;
the most significant digit is its first letter plus one. Dividing repeatedly
by $b$ recovers the word from right to left. The loop needs at most the
original code plus one iterations. Invalid numerical inputs are assigned
a failure flag; every assertion about words includes the validity test.
For $q=0$ take the empty word as the only word and omit letter operations.
This exceptional case satisfies all compatible statements below directly.

\begin{lemma}[Coding, concatenation, and length]\label{word:code}
Word validity, length, letter lookup, concatenation, prefix, suffix,
substring, and reversal are primitive recursive. Coding is injective, and
\[
 |uv|=|u|+|v|,\qquad
 \operatorname{code}(uv)=b^{|v|}\operatorname{code}(u)+
                          \operatorname{code}(v),\qquad
 0\le\operatorname{code}(w)<b^{|w|}.
\]
The last strict bound includes $w=\eps$, since $0<1$.
\end{lemma}
\begin{proof}
Division with remainder uniquely recovers the last digit and shorter prefix.
Induction on the digit count proves injectivity. The concatenation formula
follows by appending one letter at a time; the length equation follows in
the same induction. The bound is preserved by
$bc+(a+1)\le b(b^n-1)+(b-1)=b^{n+1}-1$.
To look up a letter, divide by the appropriate power of $b$ and take the
remainder. Prefixes and suffixes use division and remainder by $b^k$;
substrings are compositions of these operations. Reversal reads the letters
in decreasing position order and appends them to a fresh word. All these
loops are bounded by a supplied length or code, and their digit invariants
are bounded formulas.
\end{proof}

\begin{theorem}[Free-monoid equations and cancellation]\label{word:monoid}
Concatenation is associative, $\eps$ is its two-sided identity, and
\[
 xu=xv\Rightarrow u=v,\qquad ux=vx\Rightarrow u=v.
\]
Moreover $uv=\eps$ exactly when $u=v=\eps$.
\end{theorem}
\begin{proof}
The digit at each position of $(uv)w$ and $u(vw)$ is the same: its position
lies in the $u$, $v$, or $w$ interval in either grouping. Their lengths
agree, so injectivity proves associativity. An empty interval adds no
letters, proving the unit laws. In a cancellation equation, equality of
lengths first gives equal lengths of the uncancelled words; compare the
positions after the common prefix, or before the common suffix. Finally,
the sum of two natural lengths is zero precisely when both are zero.
\end{proof}

\begin{lemma}[Prefix comparison and Levi's lemma]\label{word:levi}
If $uv=xy$, either $|u|\le|x|$ and there is a unique $z$ with
$x=uz$ and $v=zy$, or $|x|\le|u|$ and there is a unique $z$ with
$u=xz$ and $y=zv$. At equality the two descriptions use $z=\eps$.
The applicable factorization is found by a primitive-recursive procedure.
\end{lemma}
\begin{proof}
Compare the two lengths, a decidable numerical comparison. In the first
case the first $|u|$ positions of $x$ equal $u$, since both are prefixes
of the displayed common word. Let $z$ be the remaining suffix of $x$.
Substitute $x=uz$ and cancel $u$. The other case is symmetric.
Uniqueness follows by cancellation; the suffix operation computes $z$.
\end{proof}

\begin{corollary}[Deciding elementary word equations]\label{word:equations}
The equation $ux=v$ has a solution exactly when $u$ is a prefix of $v$;
the solution is then unique and is the suffix after that prefix. The
analogous right-suffix equation is equally decidable. All factorizations
$w=uv$ are exactly the $|w|+1$ cuts at positions $0,\ldots,|w|$.
\end{corollary}
\begin{proof}
Apply the length and position equations. A solution must have length
$|v|-|u|$ and exactly the specified letters. Every cut plainly supplies
a factorization, and its left-factor length uniquely determines it.
\end{proof}

\begin{lemma}[Powers and reversal]\label{word:powers}
The power function $w^0=\eps$, $w^{n+1}=w^nw$ is primitive recursive,
and
\[
 w^{m+n}=w^mw^n,\quad |w^n|=n|w|,\quad
 (uv)^R=v^Ru^R,\quad (w^R)^R=w,\quad (w^n)^R=(w^R)^n.
\]
\end{lemma}
\begin{proof}
The first two identities are induction on the exponent with the monoid
equations. Reversal maps position $i$ to $|w|-1-i$; applying that map
twice is the identity, and its action on the two intervals of $uv$
gives the reversed concatenation. The last identity follows by induction
using this rule. Every induction compares coded words or numerical terms.
\end{proof}

\begin{theorem}[Commuting words are powers of a common word]
\label{word:commute}
For supplied words $u,v$, the equation $uv=vu$ holds exactly when there
are a word $t$ and numbers $r,s$ with $u=t^r$ and $v=t^s$.
When the equation holds, such witnesses are primitive-recursively
computable; one may take $|t|\le |u|+|v|$ and $r,s\le |u|+|v|+1$.
\end{theorem}
\begin{proof}
Common powers commute by the power equations. Conversely, if one word is
empty, use the other as $t$ and exponents $0,1$, taking $t=\eps$ if both
are empty. If both are nonempty and $|u|\le|v|$, Levi's lemma gives
$v=uz$. Cancelling the initial $u$ in $uuz=uzu$ gives $uz=zu$.
The total length of the new pair $(u,z)$ is $|v|$, strictly less than
the previous total. Apply the same procedure to this smaller pair. If
$u=t^r$ and $z=t^s$, then $v=t^{r+s}$, which reconstructs the required
answer. The other length order interchanges the words.

Implement the descent with fuel $|u|+|v|+1$ and store the substitutions
in a finite stack. Every nonterminal stage decreases the positive length
sum. The reverse pass verifies the equations for the stored pair at each
stage. For a nonempty $t$ the exponents are bounded by the corresponding
word lengths; the empty cases have the displayed bounds. Thus the proof
uses bounded descent and a quantifier-free reconstruction invariant, not
an unbounded well-founded recursion.
\end{proof}

\begin{theorem}[Counting and enumerating words]\label{word:count}
There is a primitive-recursive bijection between words of length $n$ over
$\Sigma_q$ and numbers below $q^n$. There are
$\sum_{i=0}^{N}q^i$ words of length at most $N$. If $q\ge1$, the words
$a^n$ for a fixed letter $a$ give arbitrarily many pairwise distinct words.
\end{theorem}
\begin{proof}
For $q\ge2$, read the letters themselves as base-$q$ digits, retaining
the supplied length so leading zeroes are not lost. Successive division
gives the inverse and checks both identities. For $q=1$ there is one word
of every length and one number below $1^n$. Enumerating length blocks
gives the second formula. Distinct exponents give distinct lengths in
the last assertion. This is a numerical enumeration statement; it makes
no cardinality assertion about a power set of languages.
\end{proof}

\begin{theorem}[Extension of a letter map]\label{word:hom}
A finite table $h:\Sigma_q\to\Sigma_r^*$ extends uniquely to a
concatenation-preserving map on words with $h(\eps)=\eps$.
The extension is primitive recursive. With
$M=\max_{a<q}|h(a)|$, it satisfies $|h(w)|\le M|w|$.
If every image is nonempty, then $|w|\le|h(w)|$.
\end{theorem}
\begin{proof}
Read the input from left to right and append the table entry for each
letter. Splitting this scan at a cut proves $h(uv)=h(u)h(v)$.
Any other such map must have the same value after every prefix, by
induction on its length. The length is the sum of the image lengths,
giving both bounds. The same argument proves the universal property
for a represented monoid with primitive-recursive multiplication and
proved identity and associativity laws, replacing concatenation by that
multiplication. This is a schema for each supplied represented monoid,
not a quantifier over arbitrary mathematical structures.
\end{proof}

\subsection{Primitive-recursive languages and their operations}

In this subsection $L(w)$ means $\ell(w)=1$ for a \emph{fixed}
primitive-recursive bit $\ell$, on valid words. Parameters of $\ell$ may
include numerical finite descriptions. $K=L$ means pointwise equivalence.
Put $0=\varnothing$, $1=\{\eps\}$, and $\top=\Sigma_q^*$, with the
context distinguishing the language $1$ from the number $1$.

Concatenation is the bounded predicate
\[
 (KL)(w)\ \longleftrightarrow\
 \exists i\le|w|\;
 K(w[0:i])\land L(w[i:|w|]).
\]
For nonempty $w$, membership in $L^*$ means that some choice of cuts
between its letters makes every nonempty factor belong to $L$;
there are at most $2^{|w|}$ choices. Set $L^*(\eps)=1$.
A mask of internal cuts, or an increasing sequence
$0=i_0<i_1<\cdots<i_k=|w|$, supplies the witness, with
$k\le|w|$. Empty factors can always be discarded.

\begin{theorem}[Decidable closure and witness bounds]\label{alg:prclosure}
Primitive-recursive languages are closed under union, intersection,
complement, difference, concatenation, star, reversal, and inverse word
homomorphism. A nonerasing homomorphic image of a primitive-recursive
language is primitive recursive as well.
\end{theorem}
\begin{proof}
Boolean operations use the membership bits. Concatenation checks the
$n+1$ cuts, and star checks the at most $2^n$ masks just defined, where
$n=|w|$. Reversal tests $L(w^R)$, and inverse image tests $L(h(w))$.
For the last assertion, $h(u)=w$ implies $|u|\le n$, so enumerate all
source words of length at most $n$ and test $L(u)$ and equality of the
image. These enumerations are primitive recursive even if exponentially
large. This argument deliberately excludes arbitrary erasing images:
an empty letter image removes the preimage length bound.
\end{proof}

\begin{theorem}[Boolean and concatenation algebra]\label{alg:semiring}
Union and intersection obey their usual commutative, associative,
idempotent, distributive and absorption laws. For the decidable languages
in this subsection, complement obeys both De Morgan laws and double
complement. With union as addition and concatenation as multiplication,
languages form an idempotent semiring:
\[
 (KL)M=K(LM),\quad 1L=L1=L,\quad 0L=L0=0,
\]
\[
 K(L\cup M)=KL\cup KM,\qquad (K\cup L)M=KM\cup LM.
\]
All operations here are monotone in their non-complemented arguments.
\end{theorem}
\begin{proof}
Boolean identities are checked on the finitely many possible input bits,
using their decidability. For concatenation associativity, a witness on
either side is exactly two nondecreasing cut positions
$0\le i\le j\le|w|$ giving a $K$ prefix, an $L$ middle, and an $M$
suffix. Moving parentheses retains these two cuts. Unit and zero use
the unique empty factor and the impossibility of a factor in $0$.
For distributivity, a chosen factor in a union supplies its disjunct;
conversely retain the same cut and inject that disjunct. Monotonicity
retains the witnesses and uses the corresponding membership implication.
Thus none of the identities needs excluded middle for an unbounded formula.
\end{proof}

\begin{theorem}[Star laws and least closure]\label{alg:star}
For primitive-recursive languages,
\[
 0^*=1^*=1,\quad L^*=1\cup LL^*=1\cup L^*L,\quad
 L\subseteq L^*,\quad L^*L^*=L^*,\quad (L^*)^*=L^*.
\]
The operation is monotone. It is the least represented submonoid containing
$L$: if a primitive-recursive language $X$ contains $\eps$ and $L$ and
is closed under concatenation, then $L^*\subseteq X$.
\end{theorem}
\begin{proof}
Use an empty factor list for $\eps$. A nonempty accepted word has a
first and a last nonempty factor; removing either proves the two unfoldings.
Conversely append the specified factor and remove it if empty. A word of
$L$ uses one factor unless it is empty. Joining two factor lists proves
$L^*L^*\subseteq L^*$, and multiplying by $\eps$ proves the reverse.
For $(L^*)^*$, select the inner cut mask for each outer block and flatten
the lists. Every inner block has length at most the original word length,
so the selections have a common bound $2^{|w|}$ and bounded collection
is available. Monotonicity keeps the same mask. For least closure,
multiply the factors in sequence; after $i$ factors the resulting prefix
lies in $X$. This induction uses the bit for $X$ and a supplied finite
factor table. The assertions for $0$ and $1$ follow by discarding empty
factors.
\end{proof}

\begin{proposition}[Positive closure and star identities]\label{alg:identities}
Write $L^+=LL^*$. Then
\[
 L^+=L^*L,\qquad L^*=1\cup L^+,\qquad
 (KL)^*K=K(LK)^*,\qquad
 (K\cup L)^*=K^*(LK^*)^*.
\]
Also $(1\cup L)^*=L^*$ and $L\subseteq K\Rightarrow L^*\subseteq K^*$.
\end{proposition}
\begin{proof}
For $L^+$ keep a list with at least one factor, allowing that factor to be
empty when $\eps\in L$. Separating the first or last factor proves the
first two equations, including that empty-word case. A word in
$(KL)^*K$ is a chain $k_0l_1k_1\cdots l_mk_m$ with each $k_i\in K$
and $l_i\in L$; moving the parentheses proves the sliding identity.
Empty factors cause no change and may be retained in this witness chain.
For the last displayed identity, take a factorization with each factor
tagged $K$ or $L$. Tags can be selected by their membership bits. Combine
the initial run of $K$ factors, and then every $L$ factor with the following
run of $K$ factors. This gives $K^*(LK^*)^*$. Expanding the grouped lists
gives the converse. All selections and regroupings are scans of bounded
factor tables. Adding or removing empty factors proves $(1\cup L)^*=L^*$;
the final monotonicity statement was also proved above.
\end{proof}

\begin{theorem}[Arden's lemma with an effective uniqueness proof]
\label{alg:arden}
For primitive-recursive languages $A,B$, $A^*B$ is a solution of
$X=AX\cup B$ and is contained in every primitive-recursive solution.
If $\eps\notin A$, it is the unique primitive-recursive solution.
\end{theorem}
\begin{proof}
The unfolding law gives $A^*B=B\cup A(A^*B)$. In any solution $X$,
$B\subseteq X$ and $AX\subseteq X$. Induction on the length of a
supplied list of $A$ factors proves $A^*B\subseteq X$.

Assume $\eps\notin A$ and fix $w\in X$. Test $B(w)$. If it holds,
stop. If it fails, the equation supplies $w=uv$ with $u\in A$ and
$v\in X$; choose a witnessing cut by bounded search. Since $|u|\ge1$,
the residual word is strictly shorter. Repeat. After at most $|w|$
positive cuts the residual is empty. At that stage the $AX$ alternative
is impossible, so the residual is in $B$. The stored prefixes are a
factorization in $A^*B$. The invariant records the original word as the
concatenation of the stored $A$ factors and the current residual in $X$;
it is a bounded check. This proves the reverse inclusion without an
unbounded unfolding argument. If $\eps\in A$, uniqueness need not hold:
for $A=1,B=0$, every $X$ solves the equation.
\end{proof}

\begin{theorem}[Reversal of language operations]\label{alg:reverse}
Reversal is an involution on represented languages and satisfies
\[
 (K\cup L)^R=K^R\cup L^R,\quad
 (K\cap L)^R=K^R\cap L^R,\quad
 (\overline L)^R=\overline{L^R},\quad
 (KL)^R=L^RK^R,\quad (L^*)^R=(L^R)^*.
\]
\end{theorem}
\begin{proof}
Membership in $L^R$ is tested at $w^R$ and reversal is an involution,
proving the Boolean equations. Reverse the two factors of a concatenation
and exchange their order. For star reverse every factor and the order of
the factor list; this is a bounded transformation and its own inverse.
\end{proof}

\begin{theorem}[Homomorphic image and inverse-image laws]
\label{alg:homlaws}
Inverse word homomorphism preserves union, intersection, complement, and
star inclusion $h^{-1}(L)^*\subseteq h^{-1}(L^*)$. For any finite letter
map, the existentially represented image satisfies
\[
 h(K\cup L)=h(K)\cup h(L),\quad h(KL)=h(K)h(L),\quad
 h(L^*)=h(L)^*.
\]
In the last equation an image-star witness includes a finite list of source
words. This is the witness-list presentation, with no unbounded collection
step from separate existential preimage assertions. If $h$ is nonerasing,
all these image languages have the bounded,
primitive-recursive readings of Theorem~\ref{alg:prclosure}.
\end{theorem}
\begin{proof}
The inverse-image equations simply substitute $h(w)$ into the membership
formula. If a list factors $w$ into inverse images of $L$, applying $h$
to its factors gives a factorization of $h(w)$ in $L^*$, allowing empty
images. A witness for an image is a supplied source word. For union keep
that word and its disjunct. For concatenation split the source word in the
forward direction; in the reverse direction concatenate the two supplied
source words and use $h(uv)=h(u)h(v)$.
For star apply $h$ to the source factor list, or concatenate the supplied
list of source words in the other direction. Empty image factors may be
deleted on the output side; the source list remains part of the certificate.
No unbounded collection principle is used to create such a list from
pointwise existential statements.
\end{proof}

\begin{proposition}[The source article's arithmetic syntax is regular]
\label{alg:syntax}
Let $D$ be a decimal digit, $D_+$ a nonzero digit, and put
\[
 N=0\ \cup\ D_+D^*,\qquad E=N(\mathtt{+}N)^*,\qquad
 T=E\ \cup\ E\mathtt{=}E.
\]
Then $T$ is exactly the language of numerals, sums, and one equality
described in the source article's example. In particular it has a
primitive-recursive membership test and is regular. It includes
\texttt{23+4=555} regardless of the arithmetic truth of that equation.
\end{proposition}
\begin{proof}
The numeral expression permits exactly $0$ or a nonzero first digit
followed by any digits. Splitting at plus signs gives a unique list of
nonempty numeral fields. Splitting at the optional unique equality sign
gives one or two such sums. These are exactly the permitted recursive
syntax rules: no equality may occur inside either side of the unique
equality, and each plus sign separates valid numeral or sum substrings.
Conversely, repeated splitting of any permitted sum yields those numeral
fields. This splitting decreases the substring length and has at most
the original length many steps. The expression is built from letters by
union, concatenation, and star, hence is a regular expression; the
regular-expression conversion in Section~\ref{sec:regular} supplies a
finite automaton. Evaluating decimal numerals or addition was never part
of this membership condition.
\end{proof}
